\documentclass[10pt, a4paper]{article}
\usepackage{preamble}

\title{Probability II}
\author{Luke Phillips}
\date{January 2026}

\begin{document}

\maketitle

\newpage

\tableofcontents

\newpage

\section{Probability Spaces}
There are three parts to a probability space:
the sample space $\Omega$,
this is a set
(finite or infinite).
Intuitively,
$\Omega$ consists of all outcomes of a randomised experiment.
Sigma-algebra or sigma-field $\mathcal{F}$ is a $\sigma$-field if it consists of some subsets of $\Omega$ and has the following properties:
\begin{enumerate}[label = (\alph*)]
    \item $\Omega \in \mathcal{F}$.

    \item If $A \in \mathcal{F}$,
    then $A ^ c \in \mathcal{F}$.

    \item If $A_1, A_2, A_3, \dotsc$ is a collection of subsets such that $A_n \in \mathcal{F}$ for every $n$ then $\bigcup_{n = 1}^{\infty}A_n \in \mathcal{F}$.
\end{enumerate}

\begin{definition}
    If $A \in \mathcal{F}$
    ($\sigma$-field)
    then $A$ is an event.
\end{definition}

\begin{example}
    Throw a dice:
    \[
    \Omega = \{1, 2, \dotsc, 6\}.
    \]
    \[
    \mathcal{F}_{\text{full}} = \{\text{all subsets of $\Omega$}\} \eqqcolon 2 ^ {\Omega}
    \]
    \[
    \mathcal{F}_{\text{trivial}} = \{\emptyset, \Omega\}.
    \]
\end{example}

The probability function $\P$,
$\P$ is a function from $\mathcal{F} \to [0, 1]$.
It must have two properties
\begin{enumerate}[label = (\alph*)]
    \item $\P(\Omega) = 1$.

    \item If $A_1, A_2, A_3, \dotsc$ are events
    ($A_n \in \mathcal{F}$)
    that are pairwise disjoint,
    then
    \[
    \P\left(\bigcup_{n = 1}^{\infty}A_n\right) = \sum_{n = 1}^{\infty}\P(A_n).
    \]
\end{enumerate}

Basic properties of $(\Omega, \mathcal{F}, \P)$.

\begin{proposition}
    Let $A \in \mathcal{F}$.
    Then
    \[
    \P(A ^ c) = 1 - \P(A).
    \]

    \begin{proof}
        \[
        \Omega = A \cup A ^ c
        \]
        (disjoint union)
        \[
        1 = \P(\Omega) = \P(A \cup A ^ c) = \P(A) + \P(A ^ c) \iff \P(A) = 1 - \P(A ^ c)
        \]
        using the pairwise disjoint probabilities summing and $\P(\Omega) = 1$.
    \end{proof}
\end{proposition}

\begin{proposition}
    Let $A$ and $B$ be two events
    ($A, B \in \mathcal{F}$).
    \begin{itemize}
        \item If $A \subseteq B$ then $\P(A) \leq \P(B)$.

        \item Suppose $A \subseteq B$ then $B \setminus A = B \cap A ^ c = \{\omega \in B\,:\, \omega \notin A\}$.
        Then
        \[
        \P(B) = \P(A) + \P(B \setminus A).
        \]
    \end{itemize}
    
    \begin{proof}
        \begin{itemize}
            \item
            \[
            \P(B) \geq \P(A)
            \]
            since $\P(B \setminus A) \geq 0$ in the following
            \[
            \P(B) = \P(A) + \P(B \setminus A).
            \]
            
            \item
            \[
            B = A \cup (B \setminus A)
            \]
            a disjoint union.
            Then
            \[
            \P(B) = \P(A) + \P(B \setminus A).
            \]
        \end{itemize}
    \end{proof}
\end{proposition}

\begin{proposition}[Boole's Inequality]
    Suppose $A_1, A_2, A_3, \dotsc$ are events.
    Then
    \[
    \P\left(\bigcup_{n = 1}^{\infty}A_n\right) \leq \sum_{n = 1}^{\infty}\P(A_n).
    \]

    \begin{proof}
        Let $B_1 = A_1$,
        $B_2 = A_2 \setminus B_1$,
        $B_3 = A_3 \setminus (B_1 \cup B_2)$,
        $B_n = A_n \setminus \left(\bigcup_{k = 1}^{n - 1}B_k\right)$.
        So $B_n \cap B_m = \emptyset$ when $n \neq m$ by construction.
        $B_n \subseteq A_n$ by construction.
        Observe
        \[
        \bigcup_{n = 1}^{\infty}B_n = \bigcup_{n = 1}^{\infty}A_n
        \]
        because $\bigcup_{n = 1}^{\infty}B_n \subseteq \bigcup_{n = 1}^{\infty}A_n$ since $B_n \subseteq A_n$.
        Suppose $\omega \in A_n$.
        There is a smallest index $k \leq n$ such that $\omega \in A_j$ and $\omega \notin A_1, \dotsc, A_{k - 1}$.
        Then $\omega \in B_k$.
        This implies $\omega \in \bigcup_{n = 1}^{\infty}B_n$.
        \[
        \P\left(\bigcup_{n = 1}^{\infty}A_n\right) = \P\left(\bigcup_{n = 1}^{\infty}B_n\right) = \sum_{n = 1}^{\infty}\P(B_n) \leq \sum_{n = 1}^{\infty}\P(A_n)
        \]
        since $B_n \subseteq A_n$.
    \end{proof}
\end{proposition}

\subsubsection{Generated \texorpdfstring{$\sigma$}{}-fields}

\begin{proposition}
    Suppose $\mathcal{F}_{\alpha}$,
    $\alpha \in \mathcal{I}$ is a collection of $\sigma$-fields on sample space $\Omega$.
    Then
    \[
    \mathcal{G} = \bigcap_{\alpha \in \mathcal{I}}\mathcal{F}_{\alpha}
    \]
    is also a $\sigma$-field.

    \begin{proof}
        Need to show:
        \begin{enumerate}[label = (\alph*)]
            \item $\Omega \in \mathcal{G}$.
            
            \item If $A \in \mathcal{G}$ then $A ^ c \in \mathcal{G}$.
            
            \item If $A_1, A_2, \dotsc \in \mathcal{G}$ then $\bigcup_{n = 1}^{\infty}A_n \subset \mathcal{G_n}$.
        \end{enumerate}
    
        \begin{enumerate}[label = (\alph*)]
            \item $\Omega \in \mathcal{F}_{\alpha}$ for every $\alpha \in \mathcal{I}$.
            Thus,
            $\Omega \in \bigcap_{\alpha \in \mathcal{I}}\mathcal{F}_{\alpha}$.
    
            \item $A \in \mathcal{G}$.
            Then $A \in \mathcal{F}_{\alpha}$ for every $\alpha \in \mathcal{I}$.
            Then $A ^ c \in \mathcal{F}_{\alpha}$ for every $\alpha$ then $A ^ c \in \bigcap_{\alpha \in \mathcal{I}}\mathcal{F}_{\alpha} = \mathcal{G}$.
    
            \item $A_n \in \mathcal{G}$ for every $n$.
            Thus,
            $A_n \in \mathcal{F}_{\alpha}$ for every $\alpha$ and $n$.
            Then
            \[
            \bigcup_{n = 1}^{\infty}A_n \in \mathcal{F}_{\alpha} \implies \bigcup_{n = 1}^{\infty}A_n \in \bigcap_{\alpha \in \mathcal{I}}\mathcal{F}_{\alpha} = \mathcal{G}
            \]
            for every $\alpha$.
        \end{enumerate}
        $\Omega$ is the sample space,
        $\mathcal{C} = \{A_1, \dotsc, A_k\}$ are subsets of $\Omega$.
        \[
        \sigma(\mathcal{C}) = \bigcap_{\alpha \in \mathcal{I}}\mathcal{F}_{\alpha}
        \]
        where $\mathcal{F}_{\alpha}$ ranges over all $\sigma$-field which contain every event in $\mathcal{C}$.
    \end{proof}
\end{proposition}

\begin{example}
    One event $A$.
    Find $\sigma(\{A\})$.

    Claim $\sigma(\{A\}) = \{\emptyset, \Omega, A, A ^ c\}$.
\end{example}

\begin{example}
    $\mathcal{C} = \{A, B\}$,
    $A \cap B = \emptyset$,
    $\Omega = \{A, B, (A \cup B) ^ c\}$
    \[
    \sigma(\mathcal{C}) = \{A, B, (A \cup B) ^ c, A \cup B, A \cup D, B \cup D, \underbrace{A \cup B \cup D}_{\Omega}, \emptyset\}
    \]
    $D = (A \cup B) ^ c$.
\end{example}

\begin{example}
    $\Omega = [0, 1]$.
    $\mathcal{C} = \{\text{all open sets $U \subseteq [0, 1]$}\}$.
    $\sigma(\mathcal{C})$ is called the Borel $\sigma$-field of $\Omega$.
\end{example}

\subsection{Monotone events}
\begin{definition}
    A sequence $A_n$,
    $n \geq 1$ of events is increasing if $A_n \subseteq A_{n + 1}$ for every $n$.
    
    A sequence $B_n$ is decreasing if $B_n \supseteq B_{n + 1}$ for every $n$.
\end{definition}

\begin{theorem}
    Let $A_n$ be an increasing sequence of events.
    Set $A_{\infty} = \bigcup_{n = 1}^{\infty}A_n$
    ($A_{\infty} = \lim_{n}A_n$).
    Then
    \[
    \P(A_{\infty}) = \lim_{n \to \infty}\P(A_n).
    \]
    Suppose $B_n$ are decreasing events.
    Set $B_{\infty} = \bigcap_{n = 1}^{\infty}B_n$
    ($B_{\infty} = \lim_{n}B_n$).
    Then
    \[
    \P(B_{\infty}) = \lim_{n \to \infty}\P(B_n).
    \]

    \begin{proof}
        Define $C_1 = A_1$,
        $C_2 = A_2 \setminus A_1$,
        $C_n = A_n \setminus A_{n - 1}$.
        The $C_n$ are pairwise disjoint since $A_n$ are increasing.
        Claim:
        \[
        A_{\infty} = \bigcup_{n = 1}^{\infty}C_n\quad\text{and}\quad A_n = \bigcup_{k = 1}^{n}C_k.
        \]
        $\bigcup_{n = 1}^{\infty}C_n \subseteq A_{\infty}$.
        $C_n \subseteq A_n$ so $\bigcup_{n = 1}^{\infty}C_n \subseteq \bigcup_{n = 1}^{\infty}A_n = A_{\infty}$.

        Suppose $\omega \in A_{\infty}$.
        Then $\omega \in A_n$ for some $n$.
        Let $k$ be the first index such that $\omega \in A_k$.
        Then $\omega \notin A_1, \dotsc, A_{k - 1}$.
        Then $\omega \in A_k \setminus A_{k - 1} = C_k$.
        Then $\omega \in \bigcup_{n = 1}^{\infty}C_n$ so $A_{\infty} \subseteq \bigcup_{n = 1}^{\infty}C_n$.
        \begin{align*}
            \P(A_{\infty}) &= \P\left(\bigcup_{n = 1}^{\infty}C_n\right) \\
            &= \sum_{n = 1}^{\infty}\P(C_n) \\
            &= \lim_{n \to \infty}\sum_{k = 1}^{n}\P(C_k) \\
            &= \lim_{n \to \infty}\P\left(\bigcup_{k = 1}^{n}C_k\right) \\
            &= \lim_{n \to \infty}\P(A_n).
        \end{align*}
    \end{proof}
\end{theorem}


\end{document}